\documentclass[12pt]{article}
\usepackage{amsmath,amsfonts,fullpage,amsthm,graphicx}

\newtheorem*{theorem}{Theorem}
\newtheorem*{comptheorem}{Comparison Theorem}
\newtheorem*{fact}{Fact}
\newtheorem*{goal}{Goal}
\newtheorem*{idea}{Idea}
\newtheorem*{defn}{Definition}
\newtheorem*{ex}{Example}

%Style section
%\setlength{\textheight}{10in}
%\setlength{\textwidth}{7in}
%\hoffset=-0.75in
 \pagestyle{plain}
% \setlength{\topmargin}{0in}
% \setlength{\headsep}{0in}
% \setlength{\oddsidemargin}{0in}
% \setlength{\evensidemargin}{0in}

\begin{document}


\pagestyle{empty} %use this to get rid of page numbers

\noindent
{Math 166 \hskip 1.5 truein  Names:\ \hrulefill}

\noindent
%\vskip 0.3truein
%\bigskip
%\noindent
{Improper Integrals, or the infinite fence}

\noindent  \hskip 3.5 truein  Section Number:\ \hrulefill
\medskip

\begin{goal} \rm
Figure out if your infinite fence encloses a finite or infinite amount of area!
\end{goal}
%\smallskip

\begin{enumerate}
%\setcounter{enumi}{-1}

\item Evaluate $\displaystyle\int_3^6 \displaystyle\frac{x}{\sqrt{x-3}} \ dx$ using the following steps:
\begin{enumerate}
\item Why is $\displaystyle\int_3^6 \displaystyle\frac{x}{\sqrt{x-3}} \ dx$ an improper integral?

\bigskip

\bigskip

\bigskip


\item Rewrite the improper integral as the limit of a finite integral by replacing the lower bound with a variable $a$ and taking a one-sided limit as $a \to 3^+$.

\bigskip

\bigskip

\bigskip

\bigskip

\item Evaluate the finite integral. 

\bigskip

\bigskip

\bigskip

\bigskip

\bigskip

\bigskip

\bigskip

\bigskip

\item Now take the limit.

\bigskip

\bigskip

\bigskip

\bigskip


\item State your conclusion: Does the integral converge or diverge? If it converges, what does it converge to? (That is, what is the finite area enclosed by your infinite fence?)

\bigskip

\bigskip

\bigskip

\bigskip



\end{enumerate}

\newpage

\item Use the same steps to determine for what values of $p$ the following integrals converge:
\begin{enumerate}
\item $\displaystyle\int_1^{\infty} \displaystyle\frac{1}{x^p} \ dx$.
\vfill
\item $\displaystyle\int_0^{1} \displaystyle\frac{1}{x^p} \ dx$.
\vfill
%(In this problem, you prove what we discussed in class yesterday about $p$-integrals.)

\pagebreak
\begin{comptheorem}
Suppose $0\leq f(x) \leq g(x)$ for all $x\geq a$ and $f$ and $g$ are continuous on $[a,\infty)$.
\begin{itemize}
\item
If $\int_a^{\infty} g(x) \ dx$  converges, then $\int_a^{\infty} f(x) \ dx$  converges also.
\item 
If $\int_a^{\infty} f(x) \ dx$ diverges, then $\int_a^{\infty} g(x) \ dx$  diverges also.
	\item
In any other case, we know nothing!
\end{itemize}
\end{comptheorem}
\end{enumerate}


\item Use the comparison theorem to determine whether the following improper integral converges or diverges: $\displaystyle\int_3^{\infty} \frac{5x}{(x^7+x^2)^{1/3}} \ dx$. 

\begin{enumerate}
\item What is the approximate largest term in the numerator?  What is the approximate largest term in the denominator?  (Example: for $(x^2 + 1)^{3/2}$, we would say that the largest term is approximately $(x^2)^{3/2} = x^3$.)
\vfill
\item If we delete all terms in the integral except for the two largest terms in the previous part, does the integral converge or diverge?  Why?
\vfill
\item Use the comparison theorem and part (b) to determine whether or not the original integral converges or diverges.
\vfill
\vfill
\end{enumerate}

%Recall the steps from class: 0.\ Make a guess. 1.\ Write an inequality (and check continuity). 2.\ Say why the easier integral converges/diverges. 3.\ State your conclusion.

%\bigskip
%
%\bigskip
%
%\bigskip 
%
%\bigskip
%
%\bigskip
%
%\bigskip
%
%\bigskip
%
%\bigskip
%
%\bigskip
%
%\bigskip
%
%\bigskip
%
%\bigskip
%
%\bigskip
%
%\bigskip
%
%\bigskip
%
%\bigskip
%
%\bigskip
%
%\bigskip
%
%\bigskip
\newpage
\item Use the comparison theorem to determine whether the following improper integral converges or diverges: $\displaystyle\int_1^{\infty} e^{-x^2} \ dx$.
\end{enumerate} 

\end{document}
